\subsection{Entity Linking}
\label{sec:rw:nel}
% \todo{say something about multilingual since I worked in italian}
\begin{definition}
    [Entity Linking]
    \Acf{nel} is the task of disambiguating and linking entity mentions detected in text to entities in a \acf{kr} or a \ac{kb}~\cite{entityorientedsearch2018balog,jm3}.
\end{definition}

\begin{figure}[htbp]
\centering
\begin{minipage}[t]{0.48\linewidth}
\small
\NELPerson{Giovanni Bianchi}{}{\NIL}
, residing in \NELLocation{Bologna}{https://en.wikipedia.org/wiki/Bologna}{\WikiIcon Turin} \texttt{[...]} -- plaintiff\\
and\\
\NELPerson{Davide Ricci}{}{\NIL}
, residing in \NELLocation{Milan}{https://en.wikipedia.org/wiki/Modena}{\WikiIcon Milan} \texttt{[...]} -- defendant\\
\textbf{FACTS}\\
The plaintiff summoned the defendant to court, alleging that on \NELDate{April 3, 2022} he lent him the sum of \MONEY{EUR 10,000}, as evidenced by a private written agreement signed by both parties, with the obligation to repay by \NELDate{December 31, 2022}, pursuant to \NELLaw{Article 1813 c.c.}{https://www.normattiva.it/atto/caricaDettaglioAtto?atto.dataPubblicazioneGazzetta=1942-04-04&atto.codiceRedazionale=042U0262&tipoDettaglio=multivigenza&qId=&classica=true&dataVigenza=&generaTabId=true&bloccoAggiornamentoBreadCrumb=true&title=lbl.dettaglioAtto&tabID=}{\NormattivaIcon Article 1813 c.c.} \texttt{[...]}
\subcaption{\ac{nel} on a judgment}
\label{fig:extract-nel-judgment}
\end{minipage}
\hfill
\begin{minipage}[t]{0.48\linewidth}
\small
\NELDate{[2025/04/22 -- 11:45]}%
\NELPerson{Giovanni Bianchi}{}{\NIL}:%
Meet me \NELDate{Friday, 19:00}, at %
\NELLocation{Bar Centrale}{}{\NIL}. Bring the papers.\\

\NELDate{[2025/04/22 -- 11:46]}%
\NELPerson{Davide Ricci}{}{\NIL}: Fine.\\

\NELDate{[2025/04/22 -- 18:51]}%
\NELPerson{Giovanni Bianchi}{}{\NIL}: \textit{[Voice message -- 2 min 58 sec]}\\
\quad \textit{Transcription:}\\
\quad \texttt{[...]} You spent my money on that weekend in %
\NELLocation{Milan}{https://en.wikipedia.org/wiki/Milan}{\WikiIcon Milan}, %
dinner at %
\NELLocation{Grand Hotel et de Milan}{https://en.wikipedia.org/wiki/Grand_Hotel_et_de_Milan}{\WikiIcon Grand Hotel et de Milan} %
and even those tickets for the race at the %
\NELLocation{Monza Circuit}{https://en.wikipedia.org/wiki/Monza_Circuit}{\WikiIcon Monza Circuit}. \texttt{[...]}\\
\subcaption{\ac{nel} on chat logs}
\label{fig:extract-nel-chat-logs}
\end{minipage}
\caption{\Ac{nel} application for \acl{sta} with clickable links.}
\label{fig:extracts-nel-new}
\end{figure}

An example of \ac{nel} applied the illustrative legal judgment and chat is available in \Cref{fig:extracts-nel-new}. 



While \ac{ner} identifies spans of text that correspond to named entities (\eg ``Barack Obama''), \ac{nel} goes one step further by resolving the ambiguity of the mention and associating it with a unique identifier in a \ac{kr} (\eg the Wikidata entity Barack Obama\footnote{\url{https://www.wikidata.org/entity/Q76}} or the English Wikipedia page Barack Obama\footnote{\url{https://en.wikipedia.org/wiki/Barack_Obama}}). This \acl{kg} process is crucial when mentions can correspond to multiple entities or when different surface forms all refer to the same entity.

For example, in the first case, the mention ``Lisbon'' may refer to Lisbon, Capital of Portugal\footnote{\url{https://www.wikidata.org/entity/Q597}} or to the movie Lisbon (1956 film)\footnote{\url{https://www.wikidata.org/entity/Q3203631}}. Furthermore, ambiguity can arise also for entities with the same type, such as Lisbon, Connecticut\footnote{\url{https://en.wikipedia.org/wiki/Lisbon,_Connecticut}}, or the other cities named ``Lisbon''\footnote{\url{https://en.wikipedia.org/wiki/Lisbon_(disambiguation)}}.

While, in the second case, the mentions ``Barack Obama'' and ``President Obama'' should both be linked to the entity Barack Obama, Former US President\footnote{\url{https://www.wikidata.org/entity/Q76}}.

Similarly to \ac{ner}, \acl{nel} is part of \ac{ie}, contributing to \acl{sta}, \acl{kbp}, and downstream tasks like entity-based semantic search. In \ac{sta}, besides highlighting entities identified with \ac{ner}, linking them to a \ac{kr} allows to create clickable hyperlinks that let users browse directly to entity pages (\eg Wikipedia or Wikidata).
Furthermore, in entity-based search applications, \ac{nel} enables more powerful facets and filters based on the entries in the \ac{kb} instead of on the mention surface forms.

\ac{nel} is also applied in \acl{qa} and \acf{ir}. Some \ac{qa} systems rely on entity linking to understand the entities mentioned in the question and retrieve relevant information from a \ac{kr} or a \ac{kg}~\cite{li-etal-2020-efficient,el_differences_for_qa_2022,eal_el_re_for_qa_2018}. Similarly, in \acl{ir}, linked entities support \emph{query understanding}~\cite{el_in_queries_2015}. In some cases, entity-based queries can be fully reduced to a \ac{nel} task, while in others, \ac{nel} serves as a bridge to access the relevant entity facts, as illustrated in \Cref{fig:rome-search,fig:barack-search}.

\subsubsection{Formalization}
Formally, given a set of entity mentions $M = \{m_1, m_2, \ldots\}$---for instance, detected by a \ac{ner} system---and a \acl{kr} $\mathrm{KR} = \{e_1, e_2, \ldots\}$ containing entities, a \ac{nel} system produces a mapping:
\begin{equation}
f: M \to \mathrm{KR} \cup \{\texttt{NIL}\},
\end{equation}
where \texttt{NIL} denotes that the mention does not correspond to any entry in the \acl{kr}.  
For example, considering the input sentence ``Barack Obama was born in Hawaii.'' and Wikidata as the reference \ac{kr}, the mention ``Barack Obama'' should be linked to the entity \texttt{Barack Obama} in Wikidata\footnote{\url{https://www.wikidata.org/entity/Q76}}. While in domain-specific scenarios, such as the illustrative legal judgment depicted in \Cref{fig:extract-ner-judgment}, the mention ``Giovanni Bianchi'' may not correspond to any entity in a general-purpose \ac{kb} like Wikidata, and thus should be linked to \texttt{NIL}.


\subsubsection{Literature}
Traditionally, \ac{nel} have been addressed by decomposing the task in two main steps~\cite{entity_linking_with_kb_2015}:
\begin{enumerate}
    \item \emph{Candidate generation}: retrieving a limited (\eg 100~\cite{blink}) set of possible entities from the \ac{kb} for a given mention, to filter out irrelevant entities and reduce the search space, since \ac{kr} can contain millions of entities.
    \item \emph{Candidate ranking}: ranking, or re-ranking if a preliminary ranking is available~\cite{blink}, the candidate entities to select the best match. This step can be performed with powerful but resource-intensive methods that are prohibitive to apply to the entire \ac{kb}~\cite{blink}.
\end{enumerate}

Pre-deep learning approaches for \ac{nel} relied on dictionaries, acronym expansion, or Web search engines for the candidate retrieval phase~\cite{entity_linking_with_kb_2015}.
For instance, dictionaries can be created from Wikipedia, leveraging redirect and disambiguation pages, that---respectively---can be used for entities aliases and to obtain a list of the entities sharing the same name~\cite{entity_linking_with_kb_2015}.
The expansion of acronyms served to augment the mention surface form for matching with dictionaries or search engines~\cite{entity_linking_with_kb_2015}. For example, the mention ``PSG'' can be expanded to ``Paris Saint-Germain'' to facilitate linking to the corresponding entity in a \ac{kb}. Web search engines, such as Google (filtering for Wikipedia pages only) or the Wikipedia search engine, have been used to retrieve candidate entities by querying them with the mention in the context and using the top-ranked resulting pages as candidate entities~\cite{entity_linking_with_kb_2015}.

For the ranking step, instead, both supervised and unsupervised methods have been used, either considering the mentions to link as independent or jointly disambiguating mentions in a document to exploit coherence among them~\cite{entity_linking_with_kb_2015}. These methods used a variety of features, including string similarity between the mention and entity names, entity popularity (\eg \texttt{Lisbon (1956 film)} is less popular than \texttt{Lisbon, Capital of Portugal}), the consistency of the entity types with the \ac{ner} labels, the textual context, or measures of coherence among entities linked in the same document~\cite{entity_linking_with_kb_2015}.

Some approaches treated \ac{nel} as a binary classification problem, where each candidate entity is classified as correct or incorrect for a given mention~\cite{entity_linking_with_kb_2015}. Other methods ranked candidates using learning-to-rank algorithm or by calculating mention-entity similarities using training-free vector representations~\cite{entity_linking_with_kb_2015}. To jointly link multiple entities exploiting coherence, probabilistic graphical models and graph-based methods have been employed~\cite{entity_linking_with_kb_2015}.
% \todo{better to add also citing lists of approaches from the survey}

\ac{nel} can also be treated as a multi-class classification problem, where each entity in the \ac{kb} represents a class~\cite{kar2018task}. However, the consequent large number of classes leads to suboptimal performance~\cite{Sevgili2022NeuralEL}, and this approach does not generalize to unseen entities, that have not been observed during training.

Extending \ac{nel} to handle entities unseen at training time is crucial for the \nilchallenge and incremental \ac{el}, where NIL entities are later added to the \ac{kr}. This setting is known as \emph{zero-shot} \ac{nel}~\cite{logeswaran-etal-2019-zero} and can be addressed, for example, by learning to represent new entities from brief textual descriptions~\cite{logeswaran-etal-2019-zero,blink}.

\subsubsection{Neural Entity Linking}
Neural \ac{nel} approaches typically leverage semantic vector representations of mentions and entities, which they compare to estimate a linking score~\cite{Sevgili2022NeuralEL}. For encoding the mention in the context, early neural models used techniques such as convolutional encoders~\cite{francis-landau-etal-2016-capturing,nguyen-etal-2016-joint,sorokin-gurevych-2018-mixing,sun-etal-2015-modeling} or attention between mentions and candidate entities~\cite{ganea-hofmann-2017-deep,le-titov-2019-boosting}, while more recent methods employ recurrent neural networks (RNNs)~\cite{kolitsas-etal-2018-end,gupta-etal-2017-entity,martins-etal-2019-joint} and later self-attention~\cite{logeswaran-etal-2019-zero,blink,yamada-etal-2022-global}, which has become the preferred technique for \ac{nel} encoders~\cite{Sevgili2022NeuralEL} with several approaches based on BERT~\cite{logeswaran-etal-2019-zero,blink,yamada-etal-2022-global}.

With respect to entity encoding, some methods apply techniques similar to word2vec~\cite{mikolov2013} to learn entity embeddings that lay in the same vector space as word embeddings~\cite{moreno_combining_word_entity_embeddings_2017,yamada-etal-2017-learning}. Alternatively, when considering a \ac{kg}, some approaches use \acl{kg} embeddings~\cite{pnel_pointer_network_based_end_to_end_entity_linking_2020,sorokin-gurevych-2018-mixing}, where entities and relations are embedded in a vector space preserving the \ac{kg} structure~\cite{kge_survey_2017}.
Other models, instead, propose the use of entity typing information for \ac{nel}~\cite{gupta-etal-2017-entity,khalife-vazirgiannis-2019-scalable,Raiman_Raiman_2018,Onoe_Durrett_2020,Heist2023nastylinker}, with some using the dense embeddings also at retrieval time with efficient \acf{ann} search methods~\cite{gillick-etal-2019-learning,blink}.

\subsubsection{The Bi-Encoder Architecture}

A notable architecture is the bi-encoder that from its introduction to \acl{nel}~\cite{gillick-etal-2019-learning,logeswaran-etal-2019-zero,blink}, has been employed by several studies~\cite{kassner-etal-2022-edin,Heist2023nastylinker,bhargav-etal-2022-zero,dong_reveal_the_unknown_2023,ayoola-etal-2022-refined,plekhanov2023multilingualendendentity}.

The \emph{bi-encoder} independently encodes the mention (in context) and each candidate entity (\eg using title and a brief textual description) into dense embeddings~\cite{blink,Humeau2020Poly-encoders}. A similarity measure between mention and entity can be computed using dot-product or cosine similarity over the vectors.
This allows to encode---in advance---all entities in the \ac{kr}, and---at inference time---to efficiently retrieve the most similar entities given a mention embedding~\cite{blink} via \acf{ann} search algorithms, such as HNSW~\cite{hnsw_2020}, whose retrieval times scales logarithmically with respect to the \ac{kr} size.

% On the other hand, the \emph{cross-encoder} jointly encodes the mention and a candidate entity, for instance, by concatenating mention, mention context, entity title, and entity description and passing them through an encoder-only transformer~\cite{blink,logeswaran-etal-2019-zero} to produce a mention-entity similarity score. This setup allows the model to attend over both the mention context and the entity description simultaneously, typically leading to more accurate predictions compared to bi-encoders~\cite{blink}. However, this approach must be repeated for each mention-entity pair, resulting in high computational costs that require to limit the candidate entity set, \eg to 100 candidates per mention~\cite{blink}. 

\begin{figure}[t]
    \centering
    \includegraphics[width=0.85\linewidth]{images/bi-cross-encoders-blink.png}
    \caption{Bi-encoder and cross-encoder architectures. Figure from \textcite{blink}. 
    The bi-encoder encodes mentions and entities independently, enabling fast retrieval. 
    The cross-encoder jointly encodes mention-entity pairs, allowing richer interactions at the cost of higher inference time.}
    \label{fig:bi-cross-encoders}
\end{figure}

The functioning of the bi-encoder is exemplified in \Cref{fig:bi-cross-encoders}, on the left.

Formally, the bi-encoder is composed of two networks, one for encoding mentions $T_m$, and one for encoding entities $T_e$.
This architecture is commonly referred to also as \emph{dual encoder} or \emph{two-tower} model~\cite{gillick-etal-2019-learning}.
When using a transformer like BERT~\cite{devlin-etal-2019-bert}, as in \textcite{blink}, the output consists of a vector for each input token and requires a $red(.)$ function to reduce it to a single embedding for the entire input~\cite{blink,Humeau2020Poly-encoders}. Example functions are averaging the output embeddings or considering the embedding from the \texttt{CLS} token~\cite{Humeau2020Poly-encoders}, \ie a special token whose representation from the last layer is used as an aggregate representation of the entire input sentence~\cite{devlin-etal-2019-bert}.
Thus, the mention and entity embeddings are obtained as follows.
\begin{equation}
    y_m = red(T_m(m)),
\end{equation}
\begin{equation}
    y_e = red(T_e(e)).
\end{equation}
Where the input $m$ and $e$ include, respectively, mention and entity information that can be concatenated using ad-hoc tokens. The input formats used in \textcite{blink} are
\begin{equation}
    x_m = \texttt{[CLS]} \;\text{context}_{\text{left}} \;\texttt{[M}_\texttt{s}\texttt{]} \;\text{mention} \;\texttt{[M}_\texttt{e}\texttt{]} \;\text{context}_{\text{right}} \;\texttt{[SEP]},
\end{equation}
\begin{equation}
    x_e = \texttt{[CLS]} \;\text{title} \;\texttt{[ENT]} \;\text{description} \;\texttt{[SEP]}.
\end{equation}
\texttt{CLS} and \texttt{SEP} are special BERT tokens that are always added at the beginning and at the end of the input sequence. $\;\texttt{[M}_\texttt{s}\texttt{]}$, $\;\texttt{[M}_\texttt{e}\texttt{]}$, and \texttt{[ENT]} are custom tokens for marking, in order, the start of the mention (for distinguishing mention and context tokens), the end of the mention, and the start of the entity description (dividing it from the entity title).
For additionally exploiting typing information, \textcite{Heist2023nastylinker} included, in both templates, a single token after the \texttt{CLS} representing the mention and entity types.

A first matching score from the bi-encoder can be calculated as the similarity between the two embeddings, \eg using the \emph{dot-product}:
\begin{equation}
s^{bi}_{m,e} = y_m \cdot y_e
\end{equation}

The bi-encoder networks are trained to maximize the score of the correct mention-entity pair with respect to the remaining entities in the batch~\cite{blink}. The batches can be randomly sampled or created with \emph{hard negatives}, \eg the ten highest scored entities for linking $m$ that are not the correct entity to which $m$ refers~\cite{blink}. The loss function for a mention-entity pair $(m_i, e_i)$ appears as follows~\cite{blink}:
\begin{equation}
\mathcal{L}(m_i, e_i) = -s(m_i, e_i) + \log \sum_{j=1}^{B} \exp \big( s(m_i, e_j) \big)
\end{equation}


% \paragraph{re-ranking with Cross-Encoders.}
To improve the precision of retrieved candidates, \textcite{blink} combined the bi-encoder with a \emph{cross-encoder} that jointly encodes the mention and each candidate entity in a single transformer input sequence. By concatenating the mention context, entity title, and description, the model can attend over both components simultaneously and compute a refined similarity score~\cite{blink,logeswaran-etal-2019-zero,Humeau2020Poly-encoders}.
This setup typically yields higher accuracy than the bi-encoder alone but is computationally expensive, since it must process every mention--entity pair independently. Consequently, it is used only as a \textbf{re-ranking stage} applied to the top-$k$ candidates (e.g., $k=100$) produced by the bi-encoder~\cite{blink,Heist2023nastylinker}.

% % cross
% \textcite{blink} paired the bi-encoder for retrieval and obtaining initial rankings with a \emph{cross-encoder} for more accurate re-ranking. The cross-encoder jointly encodes the mention and a candidate entity, for instance, by concatenating mention, mention context, entity title, and entity description and passing them through an encoder-only transformer~\cite{blink,logeswaran-etal-2019-zero} to produce a mention-entity similarity score. This setup allows the model to attend over both the mention context and the entity description simultaneously, typically leading to more accurate predictions compared to bi-encoders~\cite{blink}. However, this approach must be repeated for each mention-entity pair, resulting in high computational costs that require to limit the candidate entity set, \eg to 100 candidates per mention~\cite{blink}. 

% The cross-encoder, instead, can be formalized as a single transformer network $T_{ce}$ that jointly encodes the mention $m$ and the entity $e$ in the same input sequence~\cite{blink,Humeau2020Poly-encoders,logeswaran-etal-2019-zero}. 
% In this case, the input is constructed by concatenating the mention context with the entity candidate, separated by the \texttt{SEP} token, as follows:
% \begin{equation}
%     x = \texttt{[CLS]} \; x_m \; \texttt{[SEP]} \; x_e \; \texttt{[SEP]}.
% \end{equation}
% The transformer $T_{ce}$ processes the entire sequence and produces contextualized representations for all tokens, that are transformed in a single representation from the $red(.)$ function.
% \begin{equation}
%     y_{ce} = red(T_{ce}(x)).
% \end{equation}

% % Unlike the bi-encoder, the cross-encoder does not embed mentions and entities independently. 
% % Instead, all input tokens are attended to jointly, which allows the model to capture fine-grained interactions between the mention and the candidate entity description. 
% % This generally leads to higher accuracy, but at the cost of increased computational complexity, since each mention-entity pair must be encoded together during inference~\cite{blink,Humeau2020Poly-encoders}. 

% The similarity score is directly computed as a function of the cross-encoder representation, for instance using a linear projection~\cite{blink}:
% \begin{equation}
%     s^{ce}_{m,e} = W \cdot y_{ce} + b,
% \end{equation}
% where $W$ and $b$ are trainable parameters. 
% The cross-encoder is typically trained with the same objective as the bi-encoder, namely maximizing the score of the correct mention-entity pair over negatives within the batch.
% % \begin{equation}
% % \mathcal{L}^{ce}(m_i, e_i) = -s^{ce}(m_i, e_i) + \log \sum_{j=1}^{B} \exp \big( s^{ce}(m_i, e_j) \big).
% % \end{equation}


\subsubsection{Alternative Formulations of Entity Linking}

Beyond bi- and cross-encoder architectures, several works reformulate \ac{nel} as other \ac{nlp} tasks, such as \acl{qa} or text generation.
\textcite{procopio-etal-2023-entity} and \textcite{barba-etal-2022-extend}, for instance, model entity ranking as an \emph{extractive \ac{qa}} task, where the model identifies the correct entity span among a set of candidate entities concatenated with the mention context. This formulation enables processing multiple candidates in a single forward pass and allows interactions among them, though it remains computationally demanding for very large \acp{kb}.

Alternatively, \textcite{decao2021autoregressive} proposed a \emph{text-generation} formulation that performs end-to-end \ac{nel}, jointly detecting mentions and generating their linked entity names. To guarantee that the generated entities correspond to real \ac{kb} entries, constrained decoding is applied over a prefix tree built from Wikipedia titles. This approach eliminates the need to store large embedding indexes (tens of gigabytes in BLINK~\cite{blink}) and allows the model to consider previous links as contextual information.
Extending this work, the authors also addressed \emph{multilingual \acl{el}} fine-tuning a multilingual generative model~\cite{de-cao-etal-2022-multilingual}.

% These alternative formulations show that \ac{nel} can be expressed as a structured prediction or generation task, paving the way for later work combining constrained generation and large language models.


\subsubsection{End-to-End and Multilingual Entity Linking}

Recent approaches focus on efficient, end-to-end \ac{el} systems that can recognize and link entities within a single forward pass~\cite{ayoola-etal-2022-refined,plekhanov2023multilingualendendentity}. 
Unlike traditional pipelines, these models avoid separate re-ranking stages and instead learn to detect mentions and predict their linked entities jointly. 
Mention detection typically follows the BIO scheme~\cite{ramshaw-marcus-1995-text-bio-scheme}, and mention representations are derived either by averaging token embeddings~\cite{ayoola-etal-2022-refined} or using lightweight feed-forward layers~\cite{plekhanov2023multilingualendendentity}. 
Linking then proceeds as in bi-encoders, by comparing mention and entity embeddings through similarity search, often with \ac{ann}.  
ReFinED~\cite{ayoola-etal-2022-refined} also supports \emph{NIL prediction} by including NIL examples during training, while the model proposed by \textcite{plekhanov2023multilingualendendentity} extends this paradigm to a \emph{multilingual} setting.

\subsubsection{LLM-Based and Instruction-Following Approaches}

None of the previously discussed models rely on billion-parameter \acp{llm}. 
Recent work explores using them either as knowledge sources or as reasoning components for \ac{nel}.  
Some methods employ \acp{llm} to construct richer entity profiles~\cite{RYNKIEWICZ2025112185}, generating textual descriptions to improve similarity-based retrieval.  
Others integrate \acp{llm} directly into the linking pipeline.  
\textcite{xiao-etal-2023-instructed} fine-tune LLaMA-7B~\cite{touvron2023llamaopenefficientfoundation} for constrained generation of entity names, while \textcite{zhou-etal-2024-gendecider} apply LLaMA-2~\cite{touvron2023llama2openfoundation} in the re-ranking stage, selecting the correct entity---or ``None of the Candidates''---for each mention.  

Prompt-based approaches~\cite{liu-etal-2024-onenet,ding-etal-2024-chatel,geng-etal-2023-grammar} further reduce fine-tuning requirements. 
\textcite{liu-etal-2024-onenet} use in-context learning to select entities directly from candidate lists, while \textcite{ding-etal-2024-chatel} combines traditional candidate generation (via priors and BLINK~\cite{blink}) with \ac{llm}-based augmentation and final entity selection. 
\textcite{geng-etal-2023-grammar} use task-specific grammars to constrain generation, guiding \acp{llm} in \ac{nlp} tasks such as \ac{el} re-ranking without fine-tuning. 

Overall, these methods achieve higher accuracy and better domain adaptability~\cite{liu-etal-2024-onenet,ding-etal-2024-chatel,zhou-etal-2024-gendecider,xiao-etal-2023-instructed}.


However, efficiency remains a concern: among the discussed approaches, only \textcite{xiao-etal-2023-instructed} explicitly address it by reducing the number of model calls. Nonetheless, \acp{llm} are substantially larger than those of \textcite{ayoola-etal-2022-refined} and \textcite{plekhanov2023multilingualendendentity}, and cannot perform document-level \ac{nel} in a single forward pass, making their application costly for large-scale or batch annotation tasks.






% \todo{chatel has a good overview of llms for ie}
% \todo{llm for ie has a good overview of constrained generation approaches}

% \todo{recent trends, bela, rexel , refined, relik??}
% \todo{rexel at the end as multitask model}
% \todo{nil pred}
% \todo{something on clustering tooo}






